\subsection{\texorpdfstring{Application: Measure of the Euclidean ball in \(R^{n}\)}{Application: Measure of the Euclidean ball in R to the n}}

Let \(\mu\) be Lebesgue measure on \(\mathbf{R}\) , and \(\mu_n\) Lebesgue measure on \(\mathbf{R}^n\) , the product of \(n\) factors equal to \(\mu\) . We propose to calculate the measure \(V_n = \mu_n(\mathbf{B}_n)\) of the Euclidean unit ball. By Cor. 2 of Prop. 7,

\[
\mathrm{V} _ {n} = \int_ {- 1} ^ {+ 1} \mu_ {n - 1} \big (\mathbf {B} _ {n} (z _ {n}) \big) d z _ {n}.\tag{16}
\]

Now, the section \(\mathbf{B}_n(z_n)\) is the subset of \(\mathbf{R}^{n-1}\) defined by the relation \(\sum_{i=1}^{n-1} z_i^2 \leqslant 1 - z_n^2\) , in other words, it is the transform of the ball \(\mathbf{B}_{n-1}\) by the homothety with ratio \(\sqrt{1 - z_n^2}\) . But it follows immediately from Prop. 11 and the formula

\[
\alpha \int_ {- \infty} ^ {+ \infty} f (\alpha x) d x = \int_ {- \infty} ^ {+ \infty} f (z) d z
\]

for \(f \in \mathcal{K}(\mathbf{R})\) , that the image of \(\mu_{n-1}\) under the homothety \(\mathbf{x} \mapsto \alpha \mathbf{x}\) is the measure \(\alpha^{1-n} \mu_{n-1}\) . Therefore

\[
\mu_ {n - 1} \left(\mathbf {B} _ {n} \left(z _ {n}\right)\right) = \left(\sqrt {1 - z _ {n} ^ {2}}\right) ^ {n - 1} \mathrm{V} _ {n - 1}.
\]

Substitution in (16), and making the change of variable \(z_{n} = \sin \varphi\) (with \(-\frac{\pi}{2} \leqslant \varphi \leqslant \frac{\pi}{2}\) ), yields

\[
\mathrm{V} _ {n} = \mathrm{V} _ {n - 1} \int_ {- \frac {\pi}{2}} ^ {+ \frac {\pi}{2}} \cos^ {n} \varphi d \varphi = 2 \mathrm{V} _ {n - 1} \int_ {0} ^ {\frac {\pi}{2}} \cos^ {n} \varphi d \varphi .\tag{17}
\]

But (FRV, Ch. VII, §1, No.~3, formula (20))

\[
\int_ {0} ^ {\frac {\pi}{2}} \cos^ {m} \varphi d \varphi = \frac {1}{2} \frac {\Gamma \left(\frac {1}{2}\right) \Gamma \left(\frac {m + 1}{2}\right)}{\Gamma \left(\frac {m + 2}{2}\right)}
\]

and on substituting in the relation (17) and taking into account the expression for \(\Gamma(\frac{1}{2})\) (FRV, VII, §1, No.~3, formula (21)), one obtains finally

\[
\mathrm{V} _ {n} = \frac {\pi^ {n / 2}}{\Gamma \left(\frac {n}{2} + 1\right)}.\tag{18}
\]

